\documentclass[paper=a4, fontsize=11pt]{scrartcl}

\usepackage[T1]{fontenc}
\usepackage{fourier}
\usepackage[english]{babel}
\usepackage{amsmath,amsfonts,amsthm}

\usepackage{sectsty}
\allsectionsfont{\centering \normalfont\scshape}

\usepackage{fancyhdr}
\pagestyle{fancyplain}
\fancyhead{}
\fancyfoot[L]{}
\fancyfoot[C]{}
\fancyfoot[R]{\thepage}
\renewcommand{\headrulewidth}{0pt}
\renewcommand{\footrulewidth}{0pt}
\setlength{\headheight}{13.6pt}

\setlength\parindent{0pt}

% Title helpers
\newcommand{\horrule}[1]{\rule{\linewidth}{#1}}

\title{
\normalfont \normalsize
\textsc{Artificial Intelligence - 4,906,1.00 - Sheet 2 Spring 2026} \\[10pt]
\horrule{0.5pt} \\[0.4cm]
\huge Solutions \\[-2pt]
\horrule{2pt} \\[0.5cm]
}

\author{
\textbf{Group 12}\\
Valeria Pointet\\
Keanu Belo da Silva
}

\date{\normalsize\today}

\begin{document}
\maketitle

% -------------------- 1 --------------------
\section{Problem Formulation}

\subsection*{Question 1.1}
A search problem has five components.

\textbf{1. Initial state:}  
The initial state is the start position of the robot at $(0,0)$.

\textbf{2. Actions:}  
The possible actions are: Up, Down, Left, Right.  
Each action moves the robot by one cell.

\textbf{3. Transition model:}  
The transition model describes the result of each action.  
If the target cell is inside the grid and not blocked, the robot moves there.  
If the cell is blocked, the move is not allowed.

\textbf{4. Goal test:}  
The goal test checks if the robot is at position $(3,2)$.

\textbf{5. Path cost:}  
Each move has cost 1.  
So the total path cost is the number of steps.

% -------------------- 2 --------------------
\section{A*}

\subsection*{Question 2.1}
For A* search, we use:
\[
f(n) = g(n) + h(n)
\]

The current node is $(1,2)$.  
The goal node is $(4,3)$.

From the green path, we see that the robot has already taken 3 steps.  
So:
\[
g(n) = 3
\]

\textbf{Using Euclidean distance:}

\[
h(n) = \sqrt{(4-1)^2 + (3-2)^2}
\]

\[
h(n) = \sqrt{3^2 + 1^2} = \sqrt{10} \approx 3.16
\]

\[
f(n) = g(n) + h(n) = 3 + 3.16 = 6.16
\]

So with Euclidean distance:
\[
f(n) \approx 6.16
\]

\textbf{Using Manhattan distance:}

\[
h(n) = |4-1| + |3-2|
\]

\[
h(n) = 3 + 1 = 4
\]

\[
f(n) = g(n) + h(n) = 3 + 4 = 7
\]

So with Manhattan distance:
\[
f(n) = 7
\]

\subsection*{Question 2.2}
\textbf{Manhattan distance:}  
Use it when movement is only up, down, left, and right.  
It fits grid movement with four directions.

\textbf{Euclidean distance:}  
Use it when straight-line distance matters.  
It fits continuous movement or free movement in space.

\textbf{Chebyshev distance:}  
Use it when diagonal moves are allowed and have the same cost as straight moves.

\textbf{Octile distance:}  
Use it when diagonal moves are allowed, but diagonal and straight moves have different costs.

\textbf{Weighted distance:}  
Use it when some moves cost more than others.  
For example, different terrain or risk.

In this task, the robot can only move in four directions.  
So Manhattan distance is more suitable.

It matches the movement rules better than Euclidean distance.  
It also gives a tighter estimate of the real remaining path cost.  
So Manhattan distance is the more optimal heuristic in this scenario.

% -------------------- 3 --------------------
\section{Breadth-First Search}

\subsection*{Question 3.1}
Breadth-First Search expands the shallowest unexpanded node first.

Here, the robot can only move right or down.  
So BFS explores the grid level by level.

The goal is the gift cell.  
To reach it, BFS must first examine all reachable nodes in smaller depth levels.

In the best case, the goal is the first node checked on its level.  
In the worst case, the goal is the last node checked on its level.

From the grid, the goal is at depth 7 from the start.  
There are 18 reachable nodes before that level.  
On the goal level, there are 3 reachable nodes in total.

So:

\textbf{Best case:}
\[
18 + 1 = 19
\]
So the robot examines \textbf{19 distinct nodes}, excluding the start node.

\textbf{Worst case:}
\[
18 + 3 = 21
\]
So the robot examines \textbf{21 distinct nodes}, excluding the start node.

% -------------------- 4 --------------------
\section{Depth-First Search}

\subsection*{Question 4.1}
Depth-First Search expands the deepest unexpanded node first.

The worst-case time complexity of DFS is:
\[
O(b^m)
\]

Here:

\[
b = 2
\]
because this is a complete binary tree.

\[
m = 4
\]
because the terminal nodes are at depth 4.

So:
\[
O(b^m) = O(2^4) = O(16)
\]

Therefore, the worst-case time complexity is:
\[
\boxed{O(2^4) = O(16)}
\]

So the specific values are:
\[
b=2,\quad m=4
\]

This means DFS may need to explore all branches down to depth 4 in the worst case before it finds the success node.
\end{document}